\documentclass[12 pt]{article}

\textwidth=6.5in
\textheight=8.5in
\voffset=-54pt
\oddsidemargin=0in
\evensidemargin=0in
\setlength{\parskip}{8pt}
\setlength{\parindent}{0pt}
\newcommand{\tc}{\textcolor{blue}}
\usepackage[normalem]{ulem}
\usepackage[bottom]{footmisc}
\usepackage{tikz-cd}
\usepackage{amsmath, amssymb, amsthm}
\usepackage{ifthen}
\usepackage[inline]{enumitem}
\setlist{topsep=0pt}
\usepackage{xcolor}
\usepackage[pdftex]{graphicx}
\usepackage{parskip}
\usepackage{blindtext}
\usepackage{multicol}
%operator names
\DeclareMathOperator{\ker}{ker}
\DeclareMathOperator{\coker}{coker}
\DeclareMathOperator{\im}{im}
\DeclareMathOperator{\coord}{coord}
\DeclareMathOperator{\id}{id}
\DeclareMathOperator{\ob}{\mathrm{ob}}
\DeclareMathOperator{\mor}{\mathrm{mor}}
\DeclareMathOperator*{\colim}{\mathrm{colim}}
\newcommand{\op}{\mathrm{op}}
\newcommand{\co}{\mathrm{co}}
\newcommand{\Nat}{\mathrm{Nat}}
\newcommand{\Hom}{\mathrm{Hom}}
\newcommand{\Map}{\mathrm{Map}}
\newcommand{\End}{\mathrm{End}}
\newcommand{\Aut}{\mathrm{Aut}}
\newcommand{\Sym}{\mathrm{Sym}}
\newcommand{\ev}{\mathrm{ev}}
\newcommand{\Stab}{\mathrm{Stab}}
%blackboard letters
\newcommand{\FF}{\mathbb{F}}
\newcommand{\CC}{\mathbb{C}}
\newcommand{\QQ}{\mathbb{Q}}
\newcommand{\RR}{\mathbb{R}}
\newcommand{\ZZ}{\mathbb{Z}}
\newcommand{\NN}{\mathbb{N}}
\newcommand{\kk}{\mathbbe{k}}
%categories
\newcommand{\A}{\mathcal{A}}
\newcommand{\B}{\mathcal{B}}
\newcommand{\C}{\mathcal{C}}
\newcommand{\D}{\mathcal{D}}
\newcommand{\set}{\mathsf{Set}}
\newcommand{\grp}{\mathsf{Grp}}
\newcommand{\ab}{\mathsf{Ab}}
%theorems, defns, etc
\newtheorem{thm}{Theorem}[section]
\theoremstyle{definition}
\newtheorem*{defn*}{Definition}
\newtheorem*{thm*}{Theorem}
%This makes the problem environment work.
\theoremstyle{definition}
\newtheorem{problem}{Exercise}
%This makes the solution environment work.
\makeatletter\newenvironment{solution}[1][Solution]{\par\pushQED{\qed}%
\normalfont 
\topsep6\p@\@plus6\p@\relax\trivlist\item\relax{\itshape#1\@addpunct{.}}\hspace\labelsep\ignorespaces}{%
\popQED\endtrivlist\@endpefalse}
\makeatother

\newtheorem{Exinner}{Exercise}
\newenvironment{exercise}[1]{%
  \renewcommand\theExinner{#1}%
  \Exinner
}{\endExinner}

\newtheorem{manualtheoreminner}{Theorem}
\newenvironment{manualtheorem}[1]{%
  \renewcommand\themanualtheoreminner{#1}%
  \manualtheoreminner
}{\endmanualtheoreminner}

\usepackage{quiver}
\usepackage{fancyhdr}
\usepackage{graphicx}
\pagestyle{fancy}
\chead{EN.553.420.01.FA24 Probability}
\cfoot{\thepage}
    \lhead{Henna Parmar}
    \rhead{Assignment }


\begin{document}

    \begin{exercise}{1}
    \hfill

    $P(C_F) = \frac{9}{10}$
    
    $P(C_N) = \frac{1}{10}$

    $P(H_4 | C_F) = (\frac{1}{2})^2 = \frac{1}{16}$

    $P(H_4 | C_N) = 1$

    $P(H_4) = \frac{9}{10} \times \frac{1}{16} + 1 \times \frac{1}{10} = \frac{9}{160} + \frac{1}{10} = \frac{5}{32}$

    The probability the coin is double headed based on 4 heads in a row:

    $P(C_N | H_4) = \frac{P(H_4 | C_N)P(C_N)}{P(H_4)} = \frac{1 \times \frac{1}{10}}{\frac{5}{32}} = \frac{\frac{1}{10}}{\frac{5}{32}} = \frac{1}{10} \times \frac{32}{5} = \frac{32}{50} = \frac{16}{25}$

    $P(C_F | H_4) = \frac{P(H_4 | C_F)P(C_F)}{P(H_4)} = \frac{\frac{1}{16 } \times \frac{9}{10}}{\frac{5}{32}} = \frac{\frac{9}{160}}{\frac{5}{32}} = \frac{9}{160} \times \frac{32}{5} = \frac{9}{25}$

    $P(H | H_4) = P(C_F | H_4) \times \frac{1}{2} + P(C_N | H_4) \times 1$

    $P(H | H_4) = \frac{9}{25} \times \frac{1}{2} + \frac{16}{25} \times 1$

    $= \frac{9}{50} + \frac{16}{25} = \frac{41}{50}$
    \end{exercise}

    

    \begin{exercise}{2}
    \hfill

    BOOLAHUBOO

    B = 2

    O = 4

    L = 1

    A = 1

    H = 1

    U = 1

    10 letters

    $|\Omega | = \binom{10}{2, 4, 1, 1, 1, 1} = \frac{10!}{2!4!1!1!1!1!} = \frac{3,628,800}{2 \times 24 \times 1} = \frac{3,628,800}{48} = 75,600$

    $P(B) = \binom{10}{1, 4, 1, 1, 1, 1} = \frac{9!}{1!4!1!1!1!1!} = \frac{362,880}{2 \times 24 \times 1} = \frac{362,880}{48} = 15,120$

    There are 5! = 120 ways to arrange the non O letters.

    Between the 5 letters, there are 6 gaps where 1 O can be placed. 4 out of these gaps can be chosen. So,

    $= \binom{6}{4} = 15$

    Therefore,

    $P(B \cap O_B) = 120 \times 15 = 1,800$

    $P(O_B | B) = \frac{1,800}{15,120} = \frac{5}{42} \approx 0.1190$
    \end{exercise}



    \begin{exercise}{3}
    \hfill

    \textbf{(a)}
    \begin{itemize}
        \item one match has n ways of being placed
        \item for the n-1 elements there is no matches
    \end{itemize}

    $m_1(n) = r_1(n - 1) = 1 - \frac{1}{1!} + \frac{1}{2!} - \frac{1}{3!} +- ... + (-1)^{n-1} \frac{1}{(n-1)!}$

    \hfill

    \textbf{(b)}
    $p(m_1|p_1) = \frac{P(m_1 \cap p_1)}{P(p_1)} = \frac{P(m_1)}{P(p_1)} = \frac{1 - \frac{1}{1!} + \frac{1}{2!} -+ ... + \frac{(-1)^{n-1}}{(n-1)!}}{1 - \frac{1}{1!} + \frac{1}{2!} -+ ... + \frac{(-1)^{n-1}}{n!}}$

    \hfill

    \textbf{(c)} conditional prob. : $P(r_2|r_1) = 1 - P(m_1|r_1)$

    $P(r_2|r_1) = 1 - \frac{1 - \frac{1}{1!} + \frac{1}{2!} -+ ... + \frac{(-1)^{n-1}}{(n-1)!}}{1 - \frac{1}{1!} + \frac{1}{2!} -+ ... + \frac{(-1)^{n-1}}{n!}}$   
    \end{exercise}



    \begin{exercise}{4}
    \hfill

    (1, 1) (1, 2) (1, 3) (1, 4) (1, 5) (1, 6)

    (2, 1) (2, 2) (2, 3) (2, 4) (2, 5) (2, 6)

    (3, 1) (3, 2) (3, 3) (3, 4) (3, 5) (3, 6)

    (4, 1) (4, 2) (4, 3) (4, 4) (4, 5) (4, 6)

    (5, 1) (5, 2) (5, 3) (5, 4) (5, 5) (5, 6)

    (6, 1) (6, 2) (6, 3) (6, 4) (6, 5) (6, 6)

    \hfill

    $P(C_H) = \frac{1}{2}$

    $P(C_T) = \frac{1}{2}$

    $P(D_6 | C_H) = \frac{1}{6}$

    $P(D_6 | C_F) = \frac{5}{36}$

    $P(D_6) = P(D_6 | C_H) \times P(C_H) + P(D_6 | D_T) \times P(C_T)$

    $= \frac{1}{6} \times \frac{1}{2} + \frac{5}{36} \times \frac{1}{2}$

    $= \frac{1}{12} + \frac{5}{72}$

    $= \frac{11}{72}$
    \end{exercise}

    \begin{exercise}{5}
    \hfill

    6 tennis balls: 4 branch new, 2 used
    
    After someone uses a new ball it becomes a used ball.

    \textbf{(a)} Ways the roommate can grabs new balls:
    \begin{itemize}
        \item Probability I grabbed 2 new ones and roommate grabbed 2 new ones: 

        Probability 2 new ones grabbed = $P(M_n) = \frac{4}{6} \times \frac{3}{5} = \frac{2}{5}$

        Probability roommate grabs 2 new balls = $P(M_n | R_n) = \frac{2}{6} \times \frac{1}{5} = \frac{1}{15}$

        \item Probability I grabbed an old and a new one and roommate grabbed 2 new ones:

        Probability 1 new and 1 old are grabbed = $P(M_{n,o}) = \frac{2}{6} \times \frac{4}{5} = \frac{4}{15}$

        Probability roommate grabs 2 new balls = $P(M_{n,o} | R_n) = \frac{3}{6} \times \frac{2}{5} = \frac{1}{5}$

        \item Probability I grabbed a new and an old one and roommate grabbed 2 new ones

        Probability 1 old and 1 new are grabbed = $P(M_{o,n}) = \frac{4}{6} \times \frac{2}{5} = \frac{4}{15}$

        Probability roommate grabs 2 new balls = $P(M_{n,o} | R_n) = \frac{3}{6} \times \frac{2}{5} = \frac{1}{5}$

        \item Probability I grabbed 2 old ones and roommate grabbed 2 new ones:

        Probability 2 old ones grabbed = $P(M_o) = \frac{2}{6} \times \frac{1}{5} = \frac{1}{15}$

        Probability roommate grabs 2 new balls = $P(M_o | R_n) = \frac{4}{6} \times \frac{3}{5} = \frac{2}{5}$

        \item $P(R_n) = P(M_n | R_n) P(M_n) + P(M_{n,o} | R_n) P(M_{n,o}) + P(M_{o,n} | R_n) P(M_{o,n}) + P(M_o | R_n) P(M_o)$

        $P(R_n) = \frac{2}{5} \times \frac{1}{15} + \frac{4}{15} \times \frac{1}{5} + \frac{4}{15} \times \frac{1}{5} + \frac{1}{15} \times \frac{2}{5}$

        $P(R_n) = \frac{2}{75} + \frac{4}{75} + \frac{4}{75} + \frac{2}{75}$

        $P(R_n) = \frac{4}{25} = 0.16$
    \end{itemize}

    \hfill

    \textbf{(b)} Roommate grabbed 2 new balls, what is the prob. I used 2 new ones?

    $P(M_n | R_n) = \frac{(R_n | M_n) P(M_n)}{P(R_n)} = \frac{\frac{1}{15} \times \frac{2}{5}}{\frac{4}{25}} = \frac{1}{6}$
    \end{exercise}



    \begin{exercise}{6}
    \hfill

    $P(A|B) = \frac{1}{2}$

    $P(B|A) = \frac{1}{4}$

    $P(A \cup B) = \frac{1}{4}$

    \hfill

    $P(A|B) = \frac{P(A \cap B)}{P(B)} = \frac{1}{2}$
    
    $2 P(A \cap B) = P(B)$

    \hfill

    $P(B|A) = \frac{P(B \cap A)}{P(A)} = \frac{1}{4}$

    $4 P(B \cap A) = P(A) = 4 P(A \cap B)$

    \hfill

    $P(A \cup B) = P(A) + P(B) - P(A \cap B)$
    
    $\frac{3}{4} = P(A) + P(B) - P(A \cap B)$
    
    $\frac{3}{4} = 4 P(A \cap B) + 2 P(A \cap B) - P(A \cap B)$

    $\frac{3}{4} = 5 P(A \cap B)$

    $\frac{3}{20} = P(A \cap B)$

    \hfill

    $P(B) = 2 P(A \cap B)$

    $P(B) = 2 \times \frac{3}{20} = \frac{3}{10}$

    \hfill

    $P(A) = 4 P(A \cap B)$

    $P(A) = 4 \times \frac{3}{20} = \frac{3}{5}$      
    \end{exercise}



    \begin{exercise}{7}
    \hfill

    (1, 1) (1, 2) (1, 3) (1, 4) (1, 5) (1, 6)

    (2, 1) (2, 2) (2, 3) (2, 4) (2, 5) (2, 6)

    (3, 1) (3, 2) (3, 3) (3, 4) (3, 5) (3, 6)

    (4, 1) (4, 2) (4, 3) (4, 4) (4, 5) (4, 6)

    (5, 1) (5, 2) (5, 3) (5, 4) (5, 5) (5, 6)

    (6, 1) (6, 2) (6, 3) (6, 4) (6, 5) (6, 6)

    \hfill

    \textbf{8:}

    (2,6) (3,5) (4,4) (5, 3) (2, 6)

    $P(R_8) = \frac{5}{36}$

    \hfill

    \textbf{7:}

    (1, 6) (2, 5) (3, 4) (4, 3) (5, 2) (6, 1)

    $P(R_7) = \frac{6}{36} = \frac{1}{6}$

    \hfill

    \textbf{No 7 or 8 in a roll:}

    $36 - 6 - 5 = 25$

    $P(R_n) = \frac{25}{36}$

    \hfill

    $P(R_8 \cap R_7) = 0$

    $P(R_8 \cap R_7) = \frac{11}{36}$

    \hfill

    \textbf{Probability the sum of 8 happens before the sum of 7:}

    $p = \frac{P(R_8)}{P(R_8 \cap R_7)} = \frac{\frac{5}{36}}{\frac{11}{36}} = \frac{5}{11}$
    \end{exercise}



    \begin{exercise}{8}
    \hfill
    
    \begin{itemize}
        \item Fred hits a bullseye, $P(F) = 0.5$
        \item Greg hits a bullseye, $P(G) = 0.6$
        \item Fred and Greg hitting the bullseye is independent
    \end{itemize}

    \hfill

    \textbf{(a) They both hit the bullseye on a throw:}

    $P(F \cap G) = P(F)P(G) = 0.5 \times 0.6 = 0.3$

    \hfill

    \textbf{(b) Exactly one of the hits the bullseye on a throw:}

    $P(F \cap G^c) = 0.5 \times 0.4 = 0.2$

    $P(F^c \cap G) = 0.5 \times 0.6 = 0.3$

    P(exactly one hit) $= P(F \cap G^c) + P(F^c \cap G)$

    $= 0.2 + 0.3 = 0.5$

    \hfill

    \textbf{(c) Neither hits the bullseye on a throw:}

    $P(F^c \cap G^c) = 0.5 \cap 0.4 = 0.2$

    \hfill

    \textbf{(d)} M and A are independent

    $P(M) = P(F_1^c \cap F_2^c \cap F_3^c) = P(F_1^c)P(F_2^c)(P(F_3^c) = 0.5 \times 0.5 \times 0.5 = 0.125$

    $P(A) = P(G_1 \cup G_2 \cup G_3) = P(G_1) + P(G_2) + P(G_3) - P(G_1 \cap G_2) - P(G_1 \cap G_3) - P(G_2 \cap G_3) + P(G_1 \cap G_2 \cap G_3)$

    $= P(G_1) + P(G_2) + P(G_3) - P(G_1)P(G_2) - P(G_1)P(G_3) - P(G_2)P(G-3) + P(G_1)P(G_2)P(G_3)$

    $= 0.6 + 0.6 + 0.6 - 0.6 * 0.6 - 0.6 * 0.6 - 0.6 * 0.6 + 0.6 * 0.6 * 0.6$

    $= 1.8 - 0.36 - 0.36 - 0.36 + 0.216$

    $= 0.936$

    \hfill

    $P(M \cap A) = P(M)P(A) = 0.125 \times 0.936 = 0.117$   
    \end{exercise}



    \begin{exercise}{9}
    \hfill

    \textbf{(a) At least one of them hits:}

    $P(H) = 1 - P(F^c \cap G^c)$

    $P(H) = 1 - P(F^c)P(G^c)$

    $P(H) = 1 - (0.5)(0.4)$

    $P(H) = 1 - 0.2$

    $P(H) = 0.8$

    $P((F \cap G)|H) = \frac{P((F \cap G) \cap H)}{P(H)} = \frac{P(F \cap H) P(G \cap H)}{P(H)} = \frac{0.5 \times 0.6}{0.8} = \frac{0.3}{0.8} = \frac{3}{8} = 0.375$

    \hfill

    \textbf{(b)}
    $P((F^c \cap G)|H) = \frac{P((F^c \cap G) \cap H)}{P(H)} = \frac{P(F^c \cap H) P(G \cap H)}{P(H)} = \frac{0.5 \times 0.6}{0.8} = \frac{0.3}{0.8} = \frac{3}{8}  = 0.375$

    \hfill

    \textbf{(c)}
    P(exactly one hit | H) $ = \frac{P(1 hit)}{P(H)} = \frac{0.5}{0.8} = \frac{5}{8} = 0.625$
    \end{exercise}



    \begin{exercise}{10}
    \hfill

    \begin{itemize}
        \item 10 balls, 3 black and 7 white
        \item color drawn 2 more of that color put back with it
    \end{itemize}

    Prob. black balls drawn each of the first 3 trials: $\frac{3}{10} \times \frac{5}{12} \frac{7}{14} = \frac{105}{1,680} = 0.0625$
    \end{exercise}



    \begin{exercise}{11}
    \hfill

    \begin{itemize}
        \item Box A: 3 chips labeled 1-3
        \item Box B: 6 chips labeled 1-6
        \item Box C: 12 chips labeled 1-12
        \item S: Amy selects the same chip twice in a row
    \end{itemize}

    $P(A) = P(B) = P(C) = \frac{1}{3}$

    $P(S|A) = \frac{1}{3}$

    $P(S|B) = \frac{1}{6}$

    $P(S|C) = \frac{1}{12}$

    \hfill

    $P(S) = P(S|A)P(A) + P(S|B)P(B) + P(S|C)P(C)$

    $P(S) = \frac{1}{3} \times \frac{1}{3} + \frac{1}{6} \times \frac{1}{3} + \frac{1}{12} \times \frac{1}{3}$

    $P(S) = \frac{1}{9} + \frac{1}{18} + \frac{1}{36}$

    $P(S) = \frac{7}{36}$

    \hfill

    $P(B|S) = \frac{P(S|B)P(B)}{P(S)} = \frac{\frac{1}{6} \times \frac{1}{3}}{\frac{7}{36}} = \frac{2}{7}$
        
    \end{exercise}



    \begin{exercise}{12}
    \hfill

    $P(A \cap B | C) = \frac{P(A \cap B \cap C)}{P(C)}$

    \hfill

    $P(A \cap B \cap C) = P(A)P(B)P(C)$

    \hfill

    $P(A \cap B|C) = P(A)P(B)$

    \hfill

    $P(A|C) = \frac{P(A \cap C)}{P(C)} = \frac{P(A)P(C)}{P(C)} = P(A)$

    $P(B|C) = \frac{P(B \cap B)}{P(C)} = \frac{P(B)P(C)}{P(C)} = P(B)$

    \hfill

    $P(A \cap B | C) = P(A)P(B) = P(A|C)P(B|C)$

    $P(A \cap B | C) = P(A|C) P(B|C)$
    \end{exercise}
\end{document}
